\section{Design summary}

\textbf{Design decision: retain \Grind{} as a local reasoning engine and build a
structural, certificate-driven proof search system around it.} I call this
system \Forge{}. The name denotes the proposed design, not an existing released
Lean package.

A saturation engine is very effective when useful ground terms and useful
lemmas are already present. A much wider class of goals requires making one of
those things exist. For example, proving a tail-recursive reverse function
correct may require strengthening its specification from an empty accumulator
to an arbitrary accumulator. Proving an existential statement may require
constructing a function of the universally quantified inputs rather than
choosing one concrete value. Proving a quantified consequence may require
inventing an arithmetic term that does not occur in a trigger-compatible form.

\Forge{} therefore alternates between two kinds of progress:
\[
\boxed{\text{local consequence search}}
\quad\longleftrightarrow\quad
\boxed{\text{structural proof search}}.
\]
It treats \Grind{}, simplification, exact arithmetic certificates, lemma
applications, and induction as cooperating workers. Every successful worker
must return a Lean proof, a replayable certificate with a Lean soundness
argument, or explicitly unsolved subgoals. A plausible conjecture is never
entered into the proved-fact database.

The main proposed additions are a demand-driven induction and lemma-discovery
worker; a synthesis worker for source-level witnesses and arithmetic
instantiations; a nonlinear worker with exact certificates; and a controller
that preserves proof states, dependencies, and useful facts between these
operations. Higher-order and external ATP workers fit the same interface.

This is an engineering synthesis, not a claim that its constituent algorithms
are new. Aesop already supplies structured proof search; CCLemma provides
important prior art for e-graph-guided inductive lemma discovery; and existing
Lean tools provide SMT reconstruction and sum-of-squares reasoning. The value
of this proposal is the concrete integration, the demand protocol, and the
precise boundary between evidence, conjectures, and checked results.
\cite{aesop,cclemma,leansmt,sos}

\subsection{What ``more powerful'' should mean}

Here it means \emph{more goals solved automatically from comparable initial
information}, particularly goals that need structural proof steps or new
terms. It does not mean a stronger logic, a change to Lean's trusted axioms,
or a theorem that no invocation of \Grind{} could ever prove the same result.
Once the right induction hypotheses and auxiliary lemmas are supplied, a local
solver may finish a proof that it would not discover by itself.

Nor can a finite-budget heuristic promise pointwise dominance over another
heuristic while taking exactly the same resources. Two operating modes make
the tradeoff explicit. A \emph{baseline-preserving} mode runs an ordinary
\Grind{} attempt with budget $B$ and permits additional search afterward; it
preserves that attempt's successes at additional cost. A \emph{competitive}
mode shares one budget among workers; it must be benchmarked and may lose
individual baseline successes. These must not be conflated in evaluation.

\subsection{What was actually built}

The archive contains a typed first-order list prover with structural-induction
replay and accumulator generalization; a counterexample-guided polynomial
recurrence solver; an affine-witness synthesizer with Farkas certificates; a
finite-dictionary square-certificate solver; and an exact Bernstein subdivision
solver. The latter two live in one module. Their checkers are executable Python,
not formally verified Lean implementations.

The important result is qualitative as well as numerical. The list prover's
fixed-accumulator induction attempt fails, whereas a generated generalization
succeeds. The affine-witness worker solves a universally quantified strip
problem for which its constant-witness template fails. Subdivision certifies
16 of the 41 positive-polynomial cases that the root-only Bernstein check does
not certify. Those are ablations within the prototypes, not measurements
against \Grind{}.

\section{The baseline: what should not be reinvented}

The documented \Grind{} architecture already combines congruence closure,
propositional propagation, E-matching, guided case splitting, and specialized
arithmetic/algebraic reasoning. It produces ordinary Lean proof terms and uses
a shared state to let the components exchange consequences. Describing a new
e-graph with a linear arithmetic solver as a major improvement would therefore
miss the baseline.\cite{grind}

The baseline must also be dated. Lean 4.34.0 was released on September 14,
2026. Its \code{grind => bv_decide} integration encodes relevant equivalence
classes and facts learned by theory solvers and E-matching into the SAT
problem. Basic cooperation between bit-blasting and \Grind{} is consequently
\emph{existing functionality}, not a proposed \Forge{} innovation.\cite{release434}

Likewise, the linear arithmetic component handles several integer constraint
forms and already has model-based theory combination. That is not the same
thing as general model-based quantified instantiation, but it rules out a
claim that merely exchanging model-induced equalities is new.\cite{grindlia}
The release notes also describe \Grind{} discharging side conditions for
conditional \code{Sym.simp} theorems. \Forge{}'s proposal is broader
cross-worker obligation scheduling, not the invention of guarded rewriting.
\cite{release434}

\subsection{A dependency map, with integration decisions}

\begin{longtable}{@{}>{\raggedright\arraybackslash}p{0.19\textwidth}>{\raggedright\arraybackslash}p{0.37\textwidth}>{\raggedright\arraybackslash}p{0.38\textwidth}@{}}
\toprule
Component & Reuse & Integration decision \\
\midrule
\endhead
Lean core / \Grind{} & Shared equality and local-theory reasoning. & Keep as the default local closer; initially use the tactic boundary.\cite{grind,grindsource} \\
Aesop & Best-first AND--OR proof search and indexed rules. & Reuse rule-search ideas and adapters rather than claiming a new generic search tree.\cite{aesop} \\
Mathlib tactics & \code{simp}, \code{ring}, \code{linarith}, \code{nlinarith}, \code{positivity}, \code{linear_combination}. & Cheap proof reconstruction and preprocessing before expensive searches.\cite{linarith,positivity,lincomb} \\
Lean \code{sos} & Rational nonlinear certificates and a proof-producing tactic. & Optional production nonlinear backend; do not replace it with the prototype's restricted LP dictionary.\cite{sos} \\
Duper & Proof-producing automated reasoning inside Lean. & A separate lemma/first-order search worker, not an unchecked external verdict.\cite{duper} \\
\code{lean-auto} & Translation and monomorphization infrastructure. & Reuse rather than invent an unverified dependent-to-first-order erasure.\cite{leanauto} \\
\code{lean-smt} / cvc5 & SMT search with Lean reconstruction; unsupported proof steps can remain goals. & Optional backend with an explicit residual-obligation interface.\cite{leansmt} \\
Z3 & Candidate models and quantified search techniques. & Untrusted search assistance; a verdict is not a Lean proof.\cite{z3quant} \\
\code{egg} / CCLemma & Equality-saturation and inductive-lemma-search prior art. & Algorithmic references and possible isolated search processes.\cite{egg,cclemma} \\
\bottomrule
\end{longtable}

The historical \code{lean-egg} project is not selected as a new foundational
dependency: its README says it is no longer developed and recommends
\Grind{}.\cite{leanegg}

These projects do not automatically form a build-compatible package set on the
day a new Lean release appears. Pin the core toolchain and select mutually
compatible Mathlib and optional-backend revisions; keep their adapters
separate. The archive does not invent a tested \code{lake-manifest.json} for a
combination that was not built.

\section{A proof-producing obligation graph}
\status{Specified architecture. The Python prototypes implement several workers and
certificate formats, not this Lean controller.}

\subsection{Three distinct data structures}

A single global collection of ``facts'' is too ambiguous. The controller
maintains three stores with different admission rules.

The \emph{proved-fact store} contains propositions together with Lean proof
expressions valid in a particular local context. Each fact records its source,
its assumptions, and the proof dependencies needed to reconstruct it. A
worker may consume it as evidence.

The \emph{obligation graph} contains goals waiting for proofs: ordinary goals,
side conditions, candidate lemmas, induction cases, witness specifications,
and proof-reconstruction holes. An AND node is solved only after all required
children are solved. An OR node records alternative approaches to the same
goal. Failed attempts and exhausted budgets do not become logical facts.

The \emph{search store} contains models, candidate terms, speculative
rewrites, usefulness scores, and unsuccessful attempts. These may direct
search. They are not accepted as hypotheses by another proof worker. In
particular, a conjectural equality may merge nodes in an isolated exploration
graph, but must not merge nodes in the proved equality state.

\subsection{An explicit request protocol}

The following is a proposed interface, not a transcription of a released Lean
API. The goal identifier denotes an obligation in a snapshot with a known
typed context.

\begin{lstlisting}
Demand :=
  Prove(goal)
  ProveGuard(context, proposition)
  FindWitness(context, binders, specification, grammar)
  Instantiate(context, theorem, desired_consequence)
  Generalize(context, recursive_call, blocked_goal)
  FindLemma(context, left_frontier, right_frontier)
  ProveInduction(context, motive, recursor_candidates)
  Reconstruct(context, source_goal, external_certificate)

WorkerResult :=
  Proved(lean_proof, dependencies, replay_record)
  Progress(checked_facts, new_obligations, continuation)
  Candidate(untrusted_evidence, proposed_next_demands)
  Unknown(reason, diagnostics, continuation)
\end{lstlisting}

A worker returning \code{Progress} must not pretend that its residual goals
are solved. The controller may tentatively construct a parent proof containing
metavariables, but an overall success is reported only after those
metavariables are assigned and the completed expression is checked.

A counterexample belongs in \code{Candidate} unless its terms, hypotheses,
and the negation of the source conclusion have all been validated in Lean.
For an abstraction-based SMT translation, a satisfying assignment may simply
be a model of an overapproximation.

\subsection{The snapshot and its identity}

A proposed \code{GoalSnapshot} contains the environment revision, local
context, local instances, metavariable context, target expression, transparency
configuration, branch assumptions, normalization version, and worker state.
The controller must restore \emph{all} of the relevant state on backtracking.
Restoring just the visible list of hypotheses is insufficient: a worker may
have instantiated metavariables or created local declarations with identities
used by other expressions.

A persistent fact key includes its type and context identity, not just the
pretty-printed proposition. Names printed as \code{x} can denote different
free variables in sibling goals. Implicit typeclass arguments and universe
instances also matter. Two syntactically similar applications of an overloaded
operation can have different meanings because their instances differ.

Pure syntax nodes can be shared globally when their identities are stable.
Proof-bearing facts are shared only when their assumptions can be transported
into the destination context. One safe reuse method abstracts a proof over
its exact free-variable dependencies and then explicitly instantiates that
closed theorem in the destination context.

\subsection{Facts carry explanations, not only truth flags}

For example, a nonlinear worker may prove $0<c$ from polynomial constraints.
That proof justifies $c\ne0$, which unlocks a division rewrite. The rewrite's
proof yields an equality, which may enable congruence closure, which may make
an E-matching instance useful. The sequence should be represented as a proof
DAG with explicit edges, not repeated reparsing of tactic text.

A fact introduced under a split assumption $b$ has a dependency on $b$.
It cannot be used in the $\neg b$ branch. If both branches prove the same
proposition, their proofs are combined using the appropriate eliminator.
There is no ``most likely branch'' rule in the logical layer.

\section{Search control: structural changes before blind saturation}

\subsection{The main loop}

\begin{lstlisting}
solve(snapshot, budget):
  normalize_and_introduce(snapshot)
  queue = initial_demands(snapshot)
  while queue is nonempty and budget remains:
    demand = fair_priority_pop(queue)
    state = restore(demand.snapshot)
    result = select_worker(demand).run(state, local_budget)
    match result:
      Proved(p, deps, record):
        validate_in_exact_context(p, demand.goal)
        commit_proof(demand, p, deps, record)
        wake_dependent_obligations()
      Progress(facts, children, continuation):
        validate_and_commit_each_fact(facts)
        add_children_without_assuming_them(children)
        schedule(continuation)
      Candidate(evidence, suggestions):
        store_as_untrusted(evidence)
        schedule(suggestions)
      Unknown(reason, diagnostics, continuation):
        record_failure_without_logical_inference(reason)
        schedule_if_budget_policy_allows(continuation)
  return checked_root_proof_or_explicit_failure()
\end{lstlisting}

A useful first implementation is single-threaded and deterministic. Search
parallelism is optional, and is not a prerequisite for stronger reasoning.
External processes can propose certificates independently, but a single Lean
context owns proof reconstruction and commits.

\subsection{Selecting the next action}

The controller detects the shape of the obstruction rather than increasing
all budgets indiscriminately. A recursive call whose accumulator changes
suggests generalization. A goal beginning with an existential suggests witness
synthesis. A universal theorem with an arithmetic-shaped trigger suggests
constrained instantiation. A nonnegative polynomial goal with usable bounds
suggests an arithmetic certificate. An equality of functions suggests an
extensionality step before first-order reasoning.

The priority score can initially be a transparent hand-tuned combination of
estimated cost, overlap with target symbols, number of blocked dependents,
recursion relevance, and past unsuccessful attempts. These are engineering
heuristics, not calibrated probabilities. Keep a fair queue lane with growing
budgets so that a continuously regenerated cheap task cannot starve every
structural attempt. Every global budget must also account for reification,
proof reconstruction, and kernel checking.

\subsection{Do not retrieve the entire library}

Maintain indexes by conclusion head, referenced recursive definitions,
argument types, equality orientation, and the heads of premises. A blocked
obligation drives the query. Start with local hypotheses, equations of the
mentioned definitions, relevant extensionality lemmas, and a small retrieved
lemma set. Expand retrieval only if these fail.

A discrimination-tree or comparable structural index provides the first
filter; a dependency-aware score ranks the survivors. Candidate lemma
applications are elaborated in the exact goal context before they can affect
search. Text similarity or an optional language model can propose candidates,
but the system does not need either for the algorithms demonstrated here.

Aesop's best-first rule organization is a practical model for the outer search.
\Forge{} adds domain-specific demands and cross-worker continuations rather
than claiming that AND--OR search itself is novel.\cite{aesop}

\subsection{A representative cooperative proof}

Consider a program theorem that combines a tail-recursive list function with
a polynomial cost recurrence. The controller can discover an accumulator
specification, prove it by list induction, synthesize a closed form for the
cost by arithmetic induction, and pass the resulting equalities to \Grind{}
to close the original verification condition. Neither generated theorem is
asserted until its own proof is complete.

The important distinction from
\code{first | grind | aesop | nlinarith} is that unsuccessful attempts still
produce structured, reusable information: which recursive argument changed,
which equality would enable progress, which witness template was insufficient,
and which arithmetic side condition remains. This feedback directs the next
attempt, and checked intermediate results return to the shared proof state.
